(1) 取 \(m_2\) 下行为正方向，绳的加速度为 \(a=R\dot\omega\)。两个质量满足 \(m_2g-T_2=m_2a\)、\(T_1-m_1g=m_1a\)；滑轮满足 \((T_2-T_1)R-\tau_0=I a/R\)。消元得到 \(a=[(m_2-m_1)g-\tau_0/R]/[m_1+m_2+I/R^2]\)。
(2) 代回得 \(T_1=m_1(g+a)\)、\(T_2=m_2(g-a)\)，可检验 \(T_2-T_1=Ia/R^2+\tau_0/R\)。所设 \(m_2\) 下行且持续加速要求 \((m_2-m_1)g>\tau_0/R\)；若不满足，则从静止启动的问题需另外指定静摩擦规律，不能沿用恒定动摩擦力矩的运动解。
(3) 从静止走过 \(h\)，重力释放净势能 \((m_2-m_1)gh\)，轴摩擦作负功 \(-(\tau_0/R)h\)。总动能为 \([(m_1+m_2)+I/R^2]v^2/2\)，故 \(v^2=2h[(m_2-m_1)g-\tau_0/R]/[m_1+m_2+I/R^2]=2ah\)，严格与牛顿方程结果一致。
(4) 设 \(M_{\rm eff}=m_1+m_2+I/R^2\)，因为 \(\omega=v/R\)，黏滞力矩等效线性阻力 \(bv/R^2\)，故 \(M_{\rm eff}\dot v=(m_2-m_1)g-(b/R^2)v\)。初速为零时 \(v(t)=v_\infty(1-e^{-t/\tau})\)，其中 \(\tau=M_{\rm eff}R^2/b\)、\(v_\infty=(m_2-m_1)gR^2/b\)，均满足极限量纲。
